\newpage
\phantomsection
\label{prf:computational-geometry-monotone-slope-upper-hull}
\LRAProofFor{thm:computational-geometry-monotone-slope-upper-hull}

\begin{remark*}[Return]
\hyperref[thm:computational-geometry-monotone-slope-upper-hull]{Return to Theorem}
\end{remark*}
\ProofVaultURL{proof-vault/pending/computational-geometry-monotone-slope-upper-hull}

\begin{theorem*}[Monotone slope characterization of the upper hull]
Let \(P\subset\mathbb{R}^2\) be finite and in general position. Let
\[
  v_0,v_1,\dots,v_m
\]
be the vertices of the upper hull of \(P\), listed in increasing
\(x\)-coordinate, and define
\[
  s_k
  =
  \frac{y(v_{k+1})-y(v_k)}{x(v_{k+1})-x(v_k)},
  \qquad
  k=0,1,\dots,m-1.
\]
Then
\[
  s_0>s_1>\cdots>s_{m-1}.
\]
Conversely, if points \(w_0,\dots,w_m\) satisfy
\[
  x(w_0)<x(w_1)<\cdots<x(w_m)
\]
and their consecutive slopes are strictly decreasing, then every \(w_k\) is a
vertex of
\[
  \conv(\{w_0,\dots,w_m\}).
\]
The lower hull is characterized similarly, with strictly increasing consecutive
slopes.
\end{theorem*}

\begin{proof}
\textbf{Professional Standard Proof.}~\\
TODO: Prove that consecutive upper-hull edges must turn clockwise, which gives
strictly decreasing slopes, and prove the converse by showing each listed point
admits a supporting line.
\end{proof}

\begin{proof}
\textbf{Detailed Learning Proof.}~\\
TODO: Expand the analytic-geometry argument for three ordered points and show
how it applies along the whole upper chain.
\end{proof}

\begin{remark*}[Proof structure]
The proof is an ordered-slope argument: upper hull edges form a concave chain.
\end{remark*}

\begin{dependencies}
\begin{itemize}
  \item \hyperref[def:computational-geometry-convex-hull]{Convex Hull}
  \item \hyperref[def:computational-geometry-supporting-line]{Supporting Line}
  \item \hyperref[thm:computational-geometry-supporting-functional-hull-vertices]{Supporting Functional Characterization of Hull Vertices}
  \item TODO: slope of a line
\end{itemize}
\end{dependencies}

\clearpage
